\chapter{典型分布}

\section{离散分布}

\paragraph{二项分布}

$\BinDist(n,p)$：进行$n$次独立Bernoulli试验中，共成功了$X$次
\begin{equation}
	\P(X=k)=\binom nk p^k(1-p)^{n-k},\quad k=0,\ldots,n
\end{equation}

\paragraph{Pascal分布}

$\PascalDist(r,p)$：进行$X$次独立Bernoulli试验，直到成功了$r$次
\begin{equation}
	\P(X=k)=\binom{k-1}{r-1}p^r(1-p)^{k-r},\quad k=r,r+1,\ldots
\end{equation}

\paragraph{几何分布}

$\GeomDist(p)=\PascalDist(1,p)$。

\paragraph{负二项分布}

$\NegBinDist(r,p)$：进行独立Bernoulli试验，直到成功了$r$次。之前共失败了$X$次
\footnote{
	显然，Pascal分布和负二项分布只是平移的区别：$\P_{\PascalDist}(X=k+r)=\P_{\NegBinDist}(X=k)$。负二项分布名称中负的来源是
	\begin{equation}
		\binom{k+r-1}{r-1}=(-)^k\binom{-r}{k}.
	\end{equation}
}
\begin{equation}
	\P(X=k)=\binom{k+r-1}{r-1}p^r(1-p)^k,\quad k=0,1,\ldots
\end{equation}

\paragraph{超几何分布}

$\HypergeomDist(n,N,M)$：$N$个球中有$M$个有标记，从中不放回地随机抽取$n$个，其中$X$个球有标记
\begin{equation}
	\P(X=k)=\frac{\binom Mk\binom{N-M}{n-k}}{\binom N n},\quad \max(0,n+M-N)\leqslant k\leqslant\min(M,n).
\end{equation}

\paragraph{负超几何分布}

$\NegHypergeomDist(r,N,M)$：$N$个球中有$M$个有标记，从中不放回地随机抽取$X$个，
直到取到$r$个有标记的球
\begin{equation}
	\P(X=k)=\frac{\binom{k-1}{r-1}\binom{N-k}{M-r}}{\binom N M},\quad k=r,\ldots,r+N-M.
\end{equation}

\paragraph{Poisson分布}

$\PoisDist(\lambda)$：$\BinDist(n,\lambda/n)$当$n\to\infty$的极限
\begin{equation}
	\P(X=k)=\frac{\lambda^k}{k!}\e^{-\lambda},\quad k=0,1,2,\ldots
\end{equation}

\subsection{期望、方差}

\begin{center}
	\begin{tabular}{ccc}
		\toprule
		Distribution&$\E(X)$&$\Var(X)$\\
		\midrule
		$\BinDist(n,p)$&$np$&$np(1-p)$\\[1ex]
		$\PascalDist(r,p)$&$\frac rp$&$\frac rp\Bigp{\frac1p-1}$\\[2ex]
		% $\GeomDist(p)$&$\frac1p$&$\frac1p\Bigp{\frac1p-1}$\\
		$\HypergeomDist(n,N,M)$&$\frac{nM}N$&$\frac{n(N-n)M(N-M)}{N^2(N-1)}$\\[2ex]
		$\NegHypergeomDist(r,N,M)$&$\frac{r(N+1)}{M+1}$&$\frac{r(N+1)(N-M)(M-r+1)}{(M+1)^2(M+2)}$\\
		$\PoisDist(\lambda)$&$\lambda$&$\lambda$\\
		\bottomrule
	\end{tabular}
\end{center}

\paragraph{二项分布}

$X\sim\BinDist(n,p)$可以视作$X_1,\ldots,X_n$ iid $\sim\BernDist(p)$的和：$X=X_1+\cdots+X_n$，故
期望和方差
\begin{subequations}
	\begin{align}
		\E(X)&=n\E(X_i)=np,\\
		\Var(X)&=n\Var(X_i)=np(1-p).
	\end{align}
\end{subequations}
当$0<p<1$时，相邻概率之比为
\[
	\frac{\P(X=k)}{\P(X=k-1)}=\frac{n-k+1}{k}\frac p{1-p}.
\]
该比值大于1当且仅当$k<(n+1)p$。若$(n+1)p$不是整数，唯一众数是$\floor{(n+1)p}$；若$(n+1)p$为整数，则它和前一整数同为众数。退化情形$p=0,1$的众数分别为$0,n$。

% \paragraph{Pascal分布}

% 众数
% \[
% 	\frac{\P(X=k)}{\P(X=k-1)}=\frac{k+r-1}{k}(1-p)>1,
% 	\implies k<\frac{(r-1)(1-p)}{p}.
% \]
% 故
% \begin{align*}
% 	\text{Mode}=\begin{cases}
% 		r,&p>\frac{r-1}{r+1}\\
% 		(n+1)p-1\text{~or~}(n+1)p,&(n+1)p\in\NN\\
% 		\floor{(n+1)p},&(n+1)p\notin\NN
% 	\end{cases}
% \end{align*}

% \paragraph{几何分布}

% \paragraph{超几何分布}

% \paragraph{Poisson分布}

\subsection{MGF和CF}

下表记$q=1-p$。负二项分布取失败次数，其MGF要求$q\e^t<1$。

\begin{center}
	\begin{tabular}{ccc}
		\toprule
		Distribution&$\MGF_X(t)$&$\CF_X(t)$\\
		\midrule
		$\BinDist(n,p)$&$(q+p\e^t)^n$&$(q+p\e^{\i t})^n$\\
		$\NegBinDist(r,p)$&$\Bigp{\frac p{1-q\e^{t}}}^r$&$\Bigp{\frac p{1-q\e^{\i t}}}^r$\\
		$\HypergeomDist(n,N,M)$&---&---\\
		$\PoisDist(\lambda)$&$\e^{\lambda(\e^t-1)}$&$\e^{\lambda(\e^{\i t}-1)}$\\
		\bottomrule
	\end{tabular}
\end{center}

\section{连续分布}

\paragraph{均匀分布}

$\UnifDist(a,b)$
\begin{equation}
	f(x)=\frac1{b-a},\quad x\in(a,b).
\end{equation}

\paragraph{正态分布}

$\NormalDist(\mu,\sigma^2)$
\begin{equation}
	f(x)=\frac1{\sqrt{2\pi}\sigma}\exp\biggb{-\frac{(x-\mu)^2}{2\sigma^2}},\quad x\in\RR.
\end{equation}

\paragraph{指数分布}

$\ExpDist(\lambda)$
\begin{equation}
	f(x)=\lambda\e^{-\lambda x},\quad x>0.
\end{equation}

\paragraph{Gamma分布}

$\GammaDist(\alpha,\lambda)$
\begin{equation}
	f(x)=\frac{\lambda^\alpha}{\Gamma(\alpha)}x^{\alpha-1}\e^{-\lambda x},\quad x>0.
\end{equation}
其中$\alpha,\lambda>0$，Gamma函数
\[
	\Gamma(\alpha)=\intRplus x^{\alpha-1}\e^{-x}\d x,
\]
特别地，$\Gamma(x+1)=x\Gamma(x),\,\Gamma(1)=1,\,\Gamma(1/2)=\sqrt\pi,\,\Gamma(n)=(n-1)!$，且
\begin{gather*}
	\GammaDist(1,\lambda)=\ExpDist(\lambda),\,\GammaDist(n/2,1/2)=\ChiSqDist(n)\\
	f(x;\alpha,\lambda)\d x=f(\lambda x;\alpha,1)\d\lambda x
\end{gather*}

\paragraph{卡方分布}

$\ChiSqDist(n)$
\begin{equation}
	f(x)=\frac{x^{n/2-1}}{2^{n/2}\Gamma(n/2)}\e^{-x/2},\quad x>0.
\end{equation}
其中$n\in\NN_+$，特别地，$\ChiSqDist(1)$
\[
	f(x)=\frac1{\sqrt{2\pi x}}\e^{-x/2},\quad\lim_{x\to 0^+}f(x)=+\infty.
\]

\paragraph{Beta分布}

$\BetaDist(a,b)$
\begin{equation}
	f(x)=\frac1{\Beta(a,b)}x^{a-1}(1-x)^{b-1},\quad x\in(0,1).
\end{equation}
其中$a,b>0$，Beta函数
\[
	\Beta(a,b)=\int_0^1x^{a-1}(1-x)^{b-1}\d x\equiv\frac{\Gamma(a)\Gamma(b)}{\Gamma(a+b)},
\]

\paragraph{Cauchy分布}

\begin{equation}
	f(x)=\frac1\pi\frac1{1+x^2},\quad x\in\RR
\end{equation}

\paragraph{Landau分布}

略

\subsection{众数、期望、方差}

\begin{center}
	\begin{tabular}{cccc}
		\toprule
		Distribution&Mode&$\E(X)$&$\Var(X)$\\
		\midrule
		$\UnifDist(a,b)$&---&$\frac{a+b}2$&$\frac{(b-a)^2}{12}$\\
		$\NormalDist(\mu,\sigma^2)$&$\mu$&$\mu$&$\sigma^2$\\
		$\ExpDist(\lambda)$&0&$\frac1\lambda$&$\frac1{\lambda^2}$\\
		$\GammaDist(\alpha,\lambda)$&$\max(\alpha-1,0)/\lambda$&$\alpha/\lambda$&$\alpha/\lambda^2$\\
		\bottomrule
	\end{tabular}
\end{center}

指数和Gamma分布在密度集中于左边界时，将边界值0记作众数。

\section{简单随机抽样}
\label{sec:simple random sampling}

有限总体的固定数值为$Y_1,\ldots,Y_N$，均值为$\mu$，方差为$\sigma^2=N^{-1}\sum_j(Y_j-\mu)^2$。设$J_i$为不放回简单随机抽样中第$i$次抽到的编号，$X_i=Y_{J_i}$。

$A:=$样本中含有编号$j$
\[
	\P(J_i=j)=\P(J_i=j\mid A)\P(A)+0=\frac1n\cdot\frac{\binom{N-1}{n-1}}{\binom{N}{n}}=\frac1N.
\]
故
\begin{subequations}
	\begin{align}
		\E(X_i)&=\sum_{j=1}^N\P(J_i=j)Y_j=\frac1N\sum_{j=1}^NY_j=\mu;\\
		\Var(X_i)&=\sum_{j=1}^N\P(J_i=j)(Y_j-\mu)^2=\sigma^2.\\
		\E(\avg X)&=\frac1n\sum_{i=1}^n\E(X_i)=\mu.
	\end{align}
\end{subequations}
当$i\ne j$且$k\ne\ell$时，两个编号的联合抽取概率为
\[
	\P(J_i=k,J_j=\ell)=\frac1{n(n-1)}\cdot\frac{\binom{N-2}{n-2}}{\binom{N}{n}}=\frac1{N(N-1)}.
\]
故
\begin{align*}
	\E(X_iX_j)&=\frac1{N(N-1)}\sum_{k=1}^N\sum_{\ell\neq k}Y_kY_\ell=\frac1{N(N-1)}\sum_{k=1}^NY_k(N\mu-Y_k)\\
	&=\frac{N^2\mu^2-N(\sigma^2+\mu^2)}{N(N-1)}=\mu^2-\frac{\sigma^2}{N-1}.\\
	\Cov(X_i,X_j)&=\E(X_iX_j)-\E(X_i)\E(X_j)=-\frac{\sigma^2}{N-1}.
\end{align*}
故
\begin{align*}
	\Var(\avg X)&=\frac1{n^2}\biggp{\sum_{i=1}^n\Var(X_i)+\sum_{i\neq j}\Cov(X_i,X_j)}\\
	&=\frac1{n^2}\Bigp{n\sigma^2-n(n-1)\cdot\frac{\sigma^2}{N-1}}=\frac{\sigma^2}n\frac{N-n}{N-1}.\\
	\E(S^2)&=\frac n{n-1}\biggp{\frac1n\sum_{i=1}^n\E(X_i^2)-\E(\avg X^2)}\\
	&=\frac n{n-1}\biggb{\sigma^2+\mu^2-\Bigp{\frac{\sigma^2}n\frac{N-n}{N-1}+\mu^2}}=\frac{N\sigma^2}{N-1}.
\end{align*}
